% SPDX-License-Identifier: CC-BY-4.0
\documentclass[11pt,letterpaper]{article}

\usepackage[T1]{fontenc}
\usepackage{amsmath,amssymb}
\usepackage{tikz}
\usepackage{pgfplots}
\pgfplotsset{compat=1.18}
\usepackage{algorithm}
\usepackage{algpseudocode}
\usepackage{booktabs}
\usepackage{enumitem}
\usepackage[hidelinks]{hyperref}

\newcommand{\N}{\mathbb{N}}
\newcommand{\Z}{\mathbb{Z}}
\newcommand{\Q}{\mathbb{Q}}
\newcommand{\R}{\mathbb{R}}
\newcommand{\C}{\mathbb{C}}
\newcommand{\Map}[3]{#1 \colon #2 \to #3}
\newcommand{\Mod}[1]{#1\text{-}\mathrm{Mod}}
\newcommand{\GroupHom}{\mathrm{Hom}_{\mathbf{Grp}}}
\newcommand{\RingHom}{\mathrm{Hom}_{\mathbf{Ring}}}

\newtheorem{theorem}{Theorem}
\newtheorem{lemma}[theorem]{Lemma}
\newtheorem{corollary}[theorem]{Corollary}
\newtheorem{definition}[theorem]{Definition}
\newtheorem{remark}[theorem]{Remark}
\newenvironment{proof}{\noindent\textit{Proof.}}{\hfill\(\square\)}

\title{Classical Parent Joins and Terminal Words}
\author{}
\date{\today}

\begin{document}
\maketitle

\begin{abstract}
We give exact smaller-parent certificates for the classical shortcut map,
including certificates whose meeting point lies above the original seed.
Four elementary rules reduce capture to an explicitly described residual
set of natural density \(1/12\). Three further legal words reduce that set
to density \(11/216\), with uniform contraction of the parent labels.
We also give an exact valuation budget for repetitions of any fixed word
and classify terminal powers within the initial prefix \(1^a01\).
Every reduction retains both orbit clocks and the strict parent inequality.
These results do not assert capture for the remaining residual seeds.
\end{abstract}

\section{Clocks and joins}

Throughout, the domain is the positive integers and
\[
T(n)=\begin{cases}n/2,&2\mid n,\\(3n+1)/2,&2\nmid n.\end{cases}
\]
A word records the parity of the input at each shortcut tick: \(1\) is the
odd branch and \(0\) is the even branch. A legal word of length \(h\) with
\(r\) odd letters expands to \(h+r\) division-first raw ticks.
The target shortcut cycle is \(1\to2\to1\).

% Join recovery must allow the parent to have reached one before the meeting.
\begin{lemma}[Join recovery]
If \(T^b(n)=T^c(p)\) and \(p\) reaches \(1\), then \(n\) reaches \(1\).
If, in addition, \(T^d(p)=T^e(q)\), then
\[
T^{b+d}(n)=T^{c+e}(q).
\]
\end{lemma}
\begin{proof}
Choose \(k\) with \(T^k(p)=1\). If \(c\le k\), another \(k-c\) steps from
the meeting point reach \(1\). If \(c>k\), the meeting point is \(1\) or
\(2\), since the target cycle is invariant; zero or one further steps
reach \(1\). This proves the first assertion without comparing the
meeting time to the first arrival time. Applying \(T^d\) to the first
join and \(T^c\) to the second proves the displayed composition identity.
\end{proof}

Thus any rule assigning \(1\le p<n\) and a legal join permits strong
induction on the original integer label. A join alone, without this
inequality, does not furnish that induction step.

\section{An explicit residual set}

For odd \(n\), write uniquely
\[
n=2^a u-1,\qquad a=v_2(n+1)\ge1,\qquad u>0\text{ odd}.
\]
Set
\[
\mathcal R=\{2^a u-1:a\ge2,\ u>0\text{ odd},\ 3\nmid u,
                         \ 3^a u\equiv3\pmod4\}.
\]

% The four rules are ordered so their failure determines exactly the residual class.
\begin{theorem}[Four parent rules]
Every \(n>1\) outside \(\mathcal R\) has a positive parent \(p<n\) with
\(5p\le4n\) and an explicit join. Apply these rules in order:
\begin{enumerate}[label=\arabic*.]
\item If \(n\) is even, take \(p=n/2\), so \(T(n)=p\).
\item If \(n\equiv1\pmod4\), take \(p=(3n+1)/4\), so \(T^2(n)=p\).
\item If \(n\equiv2\pmod3\), take \(p=(2n-1)/3\), so \(T(p)=n\).
\item If \(a\ge2\) and \(3^a u\equiv1\pmod4\), take
\(p=2^{a-1}u-1\), so \(T^{a+2}(n)=T^{a+1}(p)\).
\end{enumerate}
\end{theorem}
\begin{proof}
The first identity and inequality are immediate. In the second case,
\(n>1\) implies \(n\ge5\); the first output is even and
\(5(3n+1)\le16n\). In the third case, the quotient is a positive odd
integer, \(3p+1=2n\), and \(p<2n/3\).

For the fourth rule the exact initial odd run is
\[
T^j(n)=3^j2^{a-j}u-1\quad(0\le j\le a).
\]
For \(j<a\) the input is odd, which proves the formula inductively;
the endpoint at \(j=a\) is even. The assumed congruence gives two even
steps from that endpoint, so
\[
T^{a+2}(n)=\frac{3^a u-1}{4}.
\]
The parent's initial run has length \(a-1\). Its endpoint is
\(3^{a-1}u-1\), congruent to \(2\) modulo \(4\). One even step and one
odd step therefore give
\[
T^{a+1}(p)=\frac{3(3^{a-1}u-1)+2}{4}
          =\frac{3^a u-1}{4}.
\]
The parent is positive because \(a\ge2\), and \(2p=n-1\).
If rules one and two fail, \(a\ge2\). If rule three also fails, then
\(3\nmid u\), since \(3\mid u\) would imply \(n\equiv2\pmod3\).
The odd integer \(3^a u\) is \(1\) or \(3\) modulo \(4\); failure of the
fourth rule leaves precisely \(\mathcal R\).
\end{proof}

\section{Three further residual certificates}

% These are coefficient identities on entire cylinders, not representative tests.
\begin{lemma}[Three cylinder identities]
For every integer \(t\ge0\),
\[
\begin{aligned}
T^3(3+8t)&=4+9t,\\
T^5(11+32t)&=10+27t,\\
T^6(7+64t)&=10+81t.
\end{aligned}
\]
The parent words in the first and third identities are \(110\) and
\(111010\); the source word in the second is \(11010\).
\end{lemma}
\begin{proof}
The first two state lists are
\[
3+8t\to5+12t\to8+18t\to4+9t
\]
and
\[
\begin{aligned}
11+32t&\to17+48t\to26+72t\\
      &\to13+36t\to20+54t\to10+27t.
\end{aligned}
\]
Every parity guard follows from the displayed constant term and even
coefficient at that input. Exact substitution in the two branches proves
every arrow. In the first list the parent satisfies
\(9(3+8t)<8(4+9t)\). In the second,
\(10+27t<11+32t\), and for \(t\ge1\) also
\(9(10+27t)\le8(11+32t)\).
The third state list is
\[
\begin{aligned}
7+64t&\to11+96t\to17+144t\to26+216t\\
     &\to13+108t\to20+162t\to10+81t.
\end{aligned}
\]
The constant parities verify \(111010\), and every coefficient is even
at its input. Substitution proves each step. Its positive parent obeys
\(9(7+64t)<8(10+81t)\), so is strictly smaller than the source.
\end{proof}

Define the smaller residual set
\[
\mathcal R_*=
\left\{n\in\mathcal R:
\begin{array}{l}
n\not\equiv4\pmod9,\quad n\not\equiv11\pmod{32},\\
n\not\equiv10\pmod{81}
\end{array}\right\}.
\]

\begin{theorem}[Residual recovery]
Every \(n>1\) outside \(\mathcal R_*\) admits a positive smaller-parent
join with \(9p\le8n\). Repeated parent reduction terminates at \(1\) or
in \(\mathcal R_*\), after at most \(\log n/\log(9/8)\) reductions.
Capture of every member of \(\mathcal R_*\) implies capture of every
positive integer. Moreover, \(d(\mathcal R_*)=11/216\).
\end{theorem}
\begin{proof}
Outside \(\mathcal R\), use the four parent rules; their factor \(4/5\)
is at most \(8/9\). Inside \(\mathcal R\), if \(n\equiv4\pmod9\), write
\(n=4+9t\) and use parent \(3+8t\) with source time zero. Otherwise, if
\(n\equiv11\pmod{32}\), write \(n=11+32t\) and use the five-tick
endpoint as parent. Here \(t=0\) would give \(u=3\), contrary to
\(3\nmid u\), so the uniform eight-ninths bound applies.
For the final uncovered class \(n\equiv10\pmod{81}\), write
\(n=10+81t\) and use parent \(7+64t\) with source time zero and
parent time six. This certificate has the same uniform bound.
All parents stay positive. The natural-number label decreases strictly,
so iteration terminates. After \(k\) reductions its positive value is at
most \((8/9)^k n\), proving the logarithmic bound. Strong induction,
using join recovery and the assumed capture on \(\mathcal R_*\), proves
the capture implication. This bounds parent depth, not forward time.

For each fixed \(a\ge2\), the residual constraints select two classes
of \(u\) modulo \(12\), and hence density \(1/(6\cdot2^a)\) in the
original \(n\)-labels. The portion with \(a\ge A\) lies in
\(n\equiv-1\pmod{2^A}\) and has upper density at most \(2^{-A}\).
Finite sums followed by this tail bound give \(d(\mathcal R)=1/12\).
Adding \(n\equiv4\pmod9\) selects one class of \(u\) modulo \(36\)
for each \(a\); the condition \(3\nmid u\) is then automatic. Its
intersection with \(\mathcal R\) has density
\[
\sum_{a\ge2}\frac1{36\cdot2^a}=\frac1{72}.
\]
The class \(n=11+32t\) has \(a=2\) and \(u=3+8t\), which always has
the required congruence modulo \(4\). Excluding \(3\mid u\) removes
exactly \(t\equiv0\pmod3\); the intersection has density \(1/48\).
Finally \(11+32t\equiv4\pmod9\) means \(t\equiv4\pmod9\), giving
intersection density \(1/288\). These two exclusions leave density
\(1/12-1/72-1/48+1/288=5/96\).
For the third class \(n\equiv10\pmod{81}\), the congruence on
\(u\) modulo \(81\) is \(2^a u\equiv11\); it automatically
excludes \(3\mid u\). Together with the required class modulo four,
this selects one class modulo \(324\) for each fixed \(a\). Its
intersection with \(\mathcal R\) therefore has density
\[
\sum_{a\ge2}\frac1{324\cdot2^a}=\frac1{648}.
\]
It is disjoint from \(n\equiv4\pmod9\), since \(10\equiv1\pmod9\).
Its intersection with \(n\equiv11\pmod{32}\) has density
\(1/(81\cdot32)=1/2592\): the two coprime congruences select one
class, which has \(a=2\), the required odd-burst congruence, and
\(3\nmid u\). The valuation tail bound above still justifies all
infinite sums. Thus the final density is
\[
\frac5{96}-\frac1{648}+\frac1{2592}=\frac{11}{216}.
\]
\end{proof}

A genuinely residual example is the entire cylinder
\[
n=31+1152t,\qquad p=27+1024t\quad(t\ge0).
\]
Here \(a=5\), \(u=1+36t\), and \(T^3(p)=n\) with \(p<n\).
The source's first five ticks are odd and strictly increasing.
This is a parent certificate before any source descent. It also shows
why an increase of the source burst length does not itself prevent a
smaller-parent join. No assertion is made that every member of
\(\mathcal R_*\) has such a certificate.

\section{Exact repeated-word budget}

For a nonempty binary word \(w\) of length \(h\), compose the rational
branches to obtain
\[
F_w(n)=\frac{Pn+B}{Q},\qquad P=3^{r(w)},\quad Q=2^h.
\]
An appended odd letter replaces \((P,Q,B)\) by
\((3P,2Q,3B+Q)\); an even letter replaces it by \((P,2Q,B)\).

% The legality criterion uses the odd suffix multiplier to recover every prefix.
\begin{lemma}[Integer cylinder criterion]
The word \(w\) is legal at an integer \(n\) if and only if \(Q\mid Pn+B\).
\end{lemma}
\begin{proof}
Legality implies integrality of the endpoint. Conversely, fix a prefix
of length \(j\) with numerator \(N_j\) and denominator \(2^j\).
The full numerator is \(3^sN_j+2^jC\), where \(s\) and \(C\) come from
the suffix. Divisibility by \(2^h\) therefore implies \(2^j\mid N_j\),
since \(3^s\) is odd. Thus every rational prefix state is integral.
For an even letter, integrality of its successor forces its input even.
For an odd letter, integrality of \((3x+1)/2\) forces its input odd.
Every letter is therefore the actual shortcut branch.
\end{proof}

\begin{theorem}[Repetition budget]
Put \(D_w(n)=(P-Q)n+B\). For every \(k\ge0\),
\[
w^k\text{ is legal at }n\quad\Longleftrightarrow\quad Q^k\mid D_w(n).
\]
If \(D_w(n)\ne0\), the maximal number of complete initial copies is
\[
\left\lfloor\frac{v_2(|D_w(n)|)}h\right\rfloor.
\]
\end{theorem}
\begin{proof}
The numerator of the \(k\)-fold rational composition is
\[
N_k=P^kn+B\sum_{j=0}^{k-1}P^{k-1-j}Q^j.
\]
Multiplication and telescoping give
\[
(P-Q)N_k=P^kD_w(n)-BQ^k.
\]
Both \(P\) and \(P-Q\) are odd, so each is a unit modulo \(Q^k\).
Consequently \(Q^k\mid N_k\) if and only if \(Q^k\mid D_w(n)\).
The cylinder criterion supplies legality. The valuation statement
follows since \(Q^k=2^{hk}\). For \(k=0\) both assertions are automatic.
\end{proof}

For \(w=110\), one has \((P,Q,B)=(9,8,5)\), hence the budget is
\(\lfloor v_2(n+5)/3\rfloor\). For a positive seed and \(P>Q\),
\(D_w(n)>0\), because \(B\ge0\); such a word cannot repeat forever.
The theorem concerns repetitions of one fixed word, not switching
between different words.

\section{Terminal powers}

% The valuation step is stated with its exact odd-length and positivity conditions.
\begin{lemma}[Divisibility at three]
For \(L>0\),
\[
3^{a+1}\mid2^L+1
\quad\Longleftrightarrow\quad
L\text{ is odd and }3^a\mid L.
\]
\end{lemma}
\begin{proof}
If \(3\mid2^L+1\), reduction modulo \(3\) makes \(L\) odd.
For odd \(L\), lifting the exponent gives
\(v_3(2^L+1)=1+v_3(L)\): apply the odd-prime sum-of-powers identity
with prime \(3\), bases \(2\) and \(1\), and odd exponent \(L\).
The prime divides \(2+1\) and does not divide either base. Comparing
the valuation to \(a+1\) proves both directions. The Lean companion
specializes this exact theorem from Mathlib \cite{Mathlib}.
\end{proof}

\begin{theorem}[Terminal-prefix classification]
Fix \(a\ge2\), let \(u>0\) be odd, and set \(n=2^a u-1\).
The legal prefix \(1^a01\) ends at a power of two if and only if
\[
u=\frac{2^{3^a t}+1}{3^{a+1}}
\quad\text{for some positive odd integer }t.
\]
Writing \(L=3^a t\), its endpoint is \(2^{L-2}\); the first hit of
\(1\) occurs at time \(a+L\). Its unique peak before that hit is
\((2^L-2)/3\), attained at time \(a\).
\end{theorem}
\begin{proof}
The odd-run formula gives \(T^a(n)=3^a u-1\), which is even.
If the following halving is odd, one more odd shortcut gives
\[
T^{a+2}(n)=\frac{3^{a+1}u-1}{4}.
\]
If this equals \(2^m\), then \(3^{a+1}u=2^{m+2}+1\).
The divisibility lemma forces \(m+2=3^a t\) with \(t>0\) odd and
recovers exactly the stated quotient for \(u\).

Conversely, the divisibility lemma makes that quotient integral. Its
numerator and denominator are odd, so \(u\) is odd, and it is positive.
Here \(L\ge9\). Reduction of \(3^{a+1}u=2^L+1\) modulo \(4\) gives
\(3^a u\equiv3\pmod4\); thus the halving after the first odd run is
odd, verifying the final letter. The endpoint is exactly \(2^{L-2}\).
It requires \(L-2\) further halvings to reach \(1\), with no earlier
hit on that tail. All states of the initial odd run increase strictly,
and its peak is \(3^a u-1=(2^L-2)/3\). The following half is smaller;
also \(3\cdot2^{L-2}<2^L-2\), so the first terminal power is smaller
than that peak. The subsequent powers decrease strictly. Every state
before the final \(1\) is greater than \(1\), proving both the unique
peak and the first-arrival assertion.
\end{proof}

\section{Formal scope}

The formal companions are in the directory
\[
\texttt{QCD.ClassicalCoalescence.001/lean/}.
\]
The file \nolinkurl{ClassicalCoalescence.lean} proves the parent-cylinder
identities above, including all branch guards and strict parent inequalities. It also proves the separate
\(79+128t\) cylinder join to \(67+108t\), whose meeting point is
\(101+162t\); that cylinder lies outside \(\mathcal R\).

In the same directory, \nolinkurl{TerminalConverse.lean} proves the
divisibility identity, the odd-run formula, the terminal
endpoint equation, and both directions of terminal-prefix classification
with the endpoint exponent retained. The first-arrival and peak
corollaries are written proofs. The complete residual recovery and
density theorem, the general repeated-word budget, and join recovery
above are also written proofs, not Lean declarations in these files.
The precise interfaces are indexed in the local library \cite{Library}.

\bibliographystyle{amsalpha}
\bibliography{ClassicalReferences}
\end{document}
